% ============================================================
\section{Introduction}
% ============================================================
The volume of Vietnamese digital text from online news, electronic commerce,
social media, and academic repositories continues to grow, making manual document
organization increasingly difficult. This study considers \emph{open-world
zero-shot multi-label text classification}, in which a system must infer the label
space and assign one or more labels to each document without task-specific labeled
data or a predefined label inventory~\cite{li2024open}. In this setting, the label
space may be expanded as additional document batches are processed. Such a capability is relevant to document collections for which a fixed taxonomy is unavailable or expensive to maintain.
Existing methods can be broadly grouped into three categories~\cite{li2024open}.
Supervised classifiers can achieve high predictive accuracy when a complete label
set and sufficient labeled examples are available, but these assumptions do not
hold in the target setting. Direct LLM-based methods can generate labels from raw
text without task-specific training, but repeated processing of long documents
increases token usage and latency and may produce entity-level or overly specific
labels. Statistical methods, including topic modeling and keyphrase extraction,
are computationally efficient but generally return topics or phrases rather than a
normalized, human-readable multi-label taxonomy. These limitations motivate a
method that combines statistical filtering with LLM-based semantic generalization.

We propose \emph{StatLM}, a pipeline comprising a statistical encoder and an LLM
decoder. The encoder removes low-centrality content and represents each document by
a compact, diverse set of keyphrases. The decoder then uses these representations
to discover, normalize, and assign coarse-grained labels. Restricting LLM processing to the final stage is intended to retain the semantic flexibility of generative models while reducing exposure to raw-text noise and limiting API usage.

The contributions of this work are as follows. First, we formulate a
hybrid encoder-decoder architecture for open-world zero-shot multi-label
classification and position it relative to supervised, statistical, and direct
LLM-based alternatives. Second, we develop a Vietnamese processing pipeline that
uses embedding-based TextRank for extractive summarization and configures KeyBERT
with an asymmetric retrieval model for keyphrase extraction. Third, we construct
and release a balanced VietnamNet corpus and evaluate the proposed pipeline at the
summarization, keyphrase-extraction, label-space, classification, and computational
levels against statistical baselines and an X-MLClass-based full-LLM baseline.

The remainder of this paper is organized as follows. Section~\ref{sec:related}
reviews the relevant methods and Vietnamese language resources.
Section~\ref{sec:method} describes the StatLM architecture, and
Section~\ref{sec:setup} presents the dataset, baselines, and evaluation protocol.
Section~\ref{sec:results} reports and discusses the experimental results.
Section~\ref{sec:conclusion} summarizes the findings, limitations, and directions
for future work.
